\documentclass[reqno]{amsart}
\usepackage[T1]{fontenc}
\usepackage{amsmath,amssymb,amsthm}
\usepackage{xr-hyper}
\usepackage[hidelinks]{hyperref}
\externaldocument[tt-]{20261125T004100--tate-twists-and-sha__workbook}
\externaldocument[lt-]{20261125T020400--local-terms-congruence-defects-and-selmer-windows__workbook}
\externaldocument[tw-]{20261126T020420--exceptional-twists-in-cyclotomic-towers__workbook}
\externaldocument[mi-]{20261126T020430--the-minus-part-by-gauss-sums__workbook}
\externaldocument[pl-]{20261126T020440--the-plus-part-by-congruences__workbook}

\newcounter{chap}
% Chapter numbers, generated from 20261125T012000--grh-workbook__guide.org
% by scripts/make-chapters.py.  Do not edit: change the order of the
% index and rerun the script (./preview-site.sh and CI do this).
%
% Usage in a sheet, after \newcounter{chap}:
%   % Chapter numbers, generated from 20261125T012000--grh-workbook__guide.org
% by scripts/make-chapters.py.  Do not edit: change the order of the
% index and rerun the script (./preview-site.sh and CI do this).
%
% Usage in a sheet, after \newcounter{chap}:
%   % Chapter numbers, generated from 20261125T012000--grh-workbook__guide.org
% by scripts/make-chapters.py.  Do not edit: change the order of the
% index and rerun the script (./preview-site.sh and CI do this).
%
% Usage in a sheet, after \newcounter{chap}:
%   \input{chapters}\chapnumber{<sheet ID>}
\makeatletter
\providecommand\chapters@set[2]{\@namedef{chapters@#1}{#2}}
\chapters@set{20261125T001000}{1}% zeros-to-ramp-sums
\chapters@set{20261125T001100}{2}% inert-primes-and-estimate-b
\chapters@set{20261125T005000}{3}% galois-fourier-and-averaging
\chapters@set{20261125T005500}{4}% columns-from-a-sphere
\chapters@set{20261125T000000}{5}% slope-method
\chapters@set{20261126T030100}{6}% a-twist-for-each-prime
\chapters@set{20261125T003000}{7}% the-chain-and-the-square-root
\chapters@set{20261125T010000}{8}% rank-lemma-and-zero-estimate
\chapters@set{20261125T010500}{9}% poincare-series-and-near-orthogonality
\chapters@set{20261124T223000}{10}% elliptic-curves-for-the-proof
\chapters@set{20261125T009500}{11}% fock-model-and-theta-intertwiners
\chapters@set{20261124T223100}{12}% complex-multiplication-and-inert-reduction
\chapters@set{20261125T008000}{13}% one-cm-determinant
\chapters@set{20261124T214500}{14}% supersingular-theta-divisibility
\chapters@set{20261124T214600}{15}% the-isotropic-plane
\chapters@set{20261125T013700}{16}% pairings-on-a-formal-group-of-height-two
\chapters@set{20261125T013500}{17}% hecke-nondegeneracy-on-the-tree
\chapters@set{20261126T030200}{18}% the-bound-is-attained
\chapters@set{20261125T009000}{19}% kronecker-limit-and-chowla-selberg
\chapters@set{20261125T002000}{20}% definite-theta-lifts
\chapters@set{20261125T015500}{21}% siegel-weil-in-the-smallest-case
\chapters@set{20261125T002100}{22}% definite-theta-lifts-and-norms
\chapters@set{20261125T016500}{23}% hecke-operators-and-theta-lifts
\chapters@set{20261125T016700}{24}% divisibility-along-the-chain
\chapters@set{20261125T004000}{25}% denominators-and-sha
\chapters@set{20261125T004100}{26}% tate-twists-and-sha
\chapters@set{20261126T020430}{27}% the-minus-part-by-gauss-sums
\chapters@set{20261125T014100}{28}% congruence-modules
\chapters@set{20261125T014500}{29}% eisenstein-congruences-and-selmer-classes
\chapters@set{20261125T014700}{30}% pseudocharacters-and-the-length-bound
\chapters@set{20261125T011000}{31}% galois-representations-black-box
\chapters@set{20261126T030300}{32}% the-data
\chapters@set{20261126T020210}{33}% bessel-poincare-series-and-their-gram-operators
\chapters@set{20261126T030400}{34}% the-seesaw-identity
\chapters@set{20261126T030500}{35}% divisibility-of-the-composite-lift
\chapters@set{20261126T030600}{36}% theorem-a-short-forms
\chapters@set{20261126T030700}{37}% the-determinant-of-a-chain
\chapters@set{20261126T030800}{38}% the-cost-of-one-lift
\chapters@set{20261126T030900}{39}% denominators-of-projections-to-lifts
\chapters@set{20261126T031000}{40}% twists-with-few-congruences
\chapters@set{20261126T031100}{41}% theorem-c-and-the-main-inequality
\providecommand\chapnumber[1]{%
  \@ifundefined{chapters@#1}%
    {\PackageError{chapters}{Sheet #1 is not listed in the index}%
      {List it in 20261125T012000--grh-workbook__guide.org\MessageBreak
       and run scripts/make-chapters.py.}}%
    {\edef\thechap{\@nameuse{chapters@#1}}}}
\makeatother
\chapnumber{<sheet ID>}
\makeatletter
\providecommand\chapters@set[2]{\@namedef{chapters@#1}{#2}}
\chapters@set{20261125T001000}{1}% zeros-to-ramp-sums
\chapters@set{20261125T001100}{2}% inert-primes-and-estimate-b
\chapters@set{20261125T005000}{3}% galois-fourier-and-averaging
\chapters@set{20261125T005500}{4}% columns-from-a-sphere
\chapters@set{20261125T000000}{5}% slope-method
\chapters@set{20261126T030100}{6}% a-twist-for-each-prime
\chapters@set{20261125T003000}{7}% the-chain-and-the-square-root
\chapters@set{20261125T010000}{8}% rank-lemma-and-zero-estimate
\chapters@set{20261125T010500}{9}% poincare-series-and-near-orthogonality
\chapters@set{20261124T223000}{10}% elliptic-curves-for-the-proof
\chapters@set{20261125T009500}{11}% fock-model-and-theta-intertwiners
\chapters@set{20261124T223100}{12}% complex-multiplication-and-inert-reduction
\chapters@set{20261125T008000}{13}% one-cm-determinant
\chapters@set{20261124T214500}{14}% supersingular-theta-divisibility
\chapters@set{20261124T214600}{15}% the-isotropic-plane
\chapters@set{20261125T013700}{16}% pairings-on-a-formal-group-of-height-two
\chapters@set{20261125T013500}{17}% hecke-nondegeneracy-on-the-tree
\chapters@set{20261126T030200}{18}% the-bound-is-attained
\chapters@set{20261125T009000}{19}% kronecker-limit-and-chowla-selberg
\chapters@set{20261125T002000}{20}% definite-theta-lifts
\chapters@set{20261125T015500}{21}% siegel-weil-in-the-smallest-case
\chapters@set{20261125T002100}{22}% definite-theta-lifts-and-norms
\chapters@set{20261125T016500}{23}% hecke-operators-and-theta-lifts
\chapters@set{20261125T016700}{24}% divisibility-along-the-chain
\chapters@set{20261125T004000}{25}% denominators-and-sha
\chapters@set{20261125T004100}{26}% tate-twists-and-sha
\chapters@set{20261126T020430}{27}% the-minus-part-by-gauss-sums
\chapters@set{20261125T014100}{28}% congruence-modules
\chapters@set{20261125T014500}{29}% eisenstein-congruences-and-selmer-classes
\chapters@set{20261125T014700}{30}% pseudocharacters-and-the-length-bound
\chapters@set{20261125T011000}{31}% galois-representations-black-box
\chapters@set{20261126T030300}{32}% the-data
\chapters@set{20261126T020210}{33}% bessel-poincare-series-and-their-gram-operators
\chapters@set{20261126T030400}{34}% the-seesaw-identity
\chapters@set{20261126T030500}{35}% divisibility-of-the-composite-lift
\chapters@set{20261126T030600}{36}% theorem-a-short-forms
\chapters@set{20261126T030700}{37}% the-determinant-of-a-chain
\chapters@set{20261126T030800}{38}% the-cost-of-one-lift
\chapters@set{20261126T030900}{39}% denominators-of-projections-to-lifts
\chapters@set{20261126T031000}{40}% twists-with-few-congruences
\chapters@set{20261126T031100}{41}% theorem-c-and-the-main-inequality
\providecommand\chapnumber[1]{%
  \@ifundefined{chapters@#1}%
    {\PackageError{chapters}{Sheet #1 is not listed in the index}%
      {List it in 20261125T012000--grh-workbook__guide.org\MessageBreak
       and run scripts/make-chapters.py.}}%
    {\edef\thechap{\@nameuse{chapters@#1}}}}
\makeatother
\chapnumber{<sheet ID>}
\makeatletter
\providecommand\chapters@set[2]{\@namedef{chapters@#1}{#2}}
\chapters@set{20261125T001000}{1}% zeros-to-ramp-sums
\chapters@set{20261125T001100}{2}% inert-primes-and-estimate-b
\chapters@set{20261125T005000}{3}% galois-fourier-and-averaging
\chapters@set{20261125T005500}{4}% columns-from-a-sphere
\chapters@set{20261125T000000}{5}% slope-method
\chapters@set{20261126T030100}{6}% a-twist-for-each-prime
\chapters@set{20261125T003000}{7}% the-chain-and-the-square-root
\chapters@set{20261125T010000}{8}% rank-lemma-and-zero-estimate
\chapters@set{20261125T010500}{9}% poincare-series-and-near-orthogonality
\chapters@set{20261124T223000}{10}% elliptic-curves-for-the-proof
\chapters@set{20261125T009500}{11}% fock-model-and-theta-intertwiners
\chapters@set{20261124T223100}{12}% complex-multiplication-and-inert-reduction
\chapters@set{20261125T008000}{13}% one-cm-determinant
\chapters@set{20261124T214500}{14}% supersingular-theta-divisibility
\chapters@set{20261124T214600}{15}% the-isotropic-plane
\chapters@set{20261125T013700}{16}% pairings-on-a-formal-group-of-height-two
\chapters@set{20261125T013500}{17}% hecke-nondegeneracy-on-the-tree
\chapters@set{20261126T030200}{18}% the-bound-is-attained
\chapters@set{20261125T009000}{19}% kronecker-limit-and-chowla-selberg
\chapters@set{20261125T002000}{20}% definite-theta-lifts
\chapters@set{20261125T015500}{21}% siegel-weil-in-the-smallest-case
\chapters@set{20261125T002100}{22}% definite-theta-lifts-and-norms
\chapters@set{20261125T016500}{23}% hecke-operators-and-theta-lifts
\chapters@set{20261125T016700}{24}% divisibility-along-the-chain
\chapters@set{20261125T004000}{25}% denominators-and-sha
\chapters@set{20261125T004100}{26}% tate-twists-and-sha
\chapters@set{20261126T020430}{27}% the-minus-part-by-gauss-sums
\chapters@set{20261125T014100}{28}% congruence-modules
\chapters@set{20261125T014500}{29}% eisenstein-congruences-and-selmer-classes
\chapters@set{20261125T014700}{30}% pseudocharacters-and-the-length-bound
\chapters@set{20261125T011000}{31}% galois-representations-black-box
\chapters@set{20261126T030300}{32}% the-data
\chapters@set{20261126T020210}{33}% bessel-poincare-series-and-their-gram-operators
\chapters@set{20261126T030400}{34}% the-seesaw-identity
\chapters@set{20261126T030500}{35}% divisibility-of-the-composite-lift
\chapters@set{20261126T030600}{36}% theorem-a-short-forms
\chapters@set{20261126T030700}{37}% the-determinant-of-a-chain
\chapters@set{20261126T030800}{38}% the-cost-of-one-lift
\chapters@set{20261126T030900}{39}% denominators-of-projections-to-lifts
\chapters@set{20261126T031000}{40}% twists-with-few-congruences
\chapters@set{20261126T031100}{41}% theorem-c-and-the-main-inequality
\providecommand\chapnumber[1]{%
  \@ifundefined{chapters@#1}%
    {\PackageError{chapters}{Sheet #1 is not listed in the index}%
      {List it in 20261125T012000--grh-workbook__guide.org\MessageBreak
       and run scripts/make-chapters.py.}}%
    {\edef\thechap{\@nameuse{chapters@#1}}}}
\makeatother

\ifcsname chapters@20261126T020450\endcsname
  \chapnumber{20261126T020450}
\else
  \renewcommand{\thechap}{?}
\fi
\renewcommand{\thesection}{\thechap.\arabic{section}}
\theoremstyle{definition}
\newtheorem{problem}{Problem}[section]
\theoremstyle{plain}
\newtheorem{theorem}{Theorem}[section]
\newtheorem{lemma}[theorem]{Lemma}
\theoremstyle{remark}
\newtheorem*{leanremark}{In the Lean formalization}

\newcommand{\Z}{\mathbb{Z}}
\newcommand{\Q}{\mathbb{Q}}
\newcommand{\Zl}{\Z_\ell}
\newcommand{\Bc}{\mathcal{B}}
\newcommand{\Ec}{\mathcal{E}}
\newcommand{\Gal}{\operatorname{Gal}}
\newcommand{\Sha}{\mathrm{Sha}}
\newcommand{\tl}[1]{\texttt{#1}}
\setlength{\emergencystretch}{3em}

\title{Windows of scales}
\author{}
\date{November 26, 2026}

\begin{document}
\maketitle
\refstepcounter{chap}\label{chap}

\paragraph{Prerequisites.}
Chapter~\ref{lt-chap}, Section~\ref{lt-sec:window}: window pairs, Theorem~\ref{lt-hyp:window}, Theorem~\ref{lt-thm:signs} and Problem~\ref{lt-pr:conductor}; there the page exponent is written \(e\) and the twist \(i\).  Chapter~\ref{tt-chap}: the lengths \(h_{\ell,e}(F)\) and the sum \(V(E,t)\) of Section~\ref{tt-sec:budget}.  Chapter~\ref{tw-chap}: layers, Lemma~\ref{tw-lem:tower}, Theorems~\ref{tw-thm:fw}, \ref{tw-thm:H2} and~\ref{tw-thm:point}, with the constant \(C_H\), and Problems~\ref{tw-pr:layer-disc} and~\ref{tw-pr:H2-small}.  Chapter~\ref{mi-chap}: the fields \(E(\psi,a,q,d)\), the two signs, and Theorem~\ref{mi-thm:minus}.  Chapter~\ref{pl-chap}: the bad primes \(\Bc_\psi\) and Theorem~\ref{pl-thm:plus}.

Throughout, \(\psi\) is a primitive Dirichlet character of conductor \(N_\psi\) and order \(o\), \(0<\varepsilon<1/52\), \(b=\exp(x^\varepsilon)\), and \(a(x)\) is a \emph{page exponent}: an integer with \(x^{2+8\varepsilon}\le2^{a(x)}\le2x^{2+8\varepsilon}\) for \(x\ge1\).  We write \(E_x=E(\psi,a(x),q,d)\), and \(N_x\) for its modulus \(\operatorname{lcm}(2^{a(x)+2},N_\psi)qd\).

\section{The statement}\label{sec:statement}

Chapter~\ref{lt-chap} states the bound that the argument needs as Theorem~\ref{lt-hyp:window}; in \cite{Adr} it is a hypothesis.  In the Lean formalization it is a theorem, and this chapter outlines its proof.

\begin{theorem}[Selmer bound on windows]\label{thm:vwin}
For \(\psi\), \(\varepsilon\) and \(a\) as above there are \(J\ge1\) and \(X_0\) such that for every \(X\ge X_0\) and every window pair \((q,d)\) at \(X^{1/J}\) there is \(Y\) with \(X^{1/J}\le Y\) and \(Y^2\le X\) such that, for every \(x\in[Y,Y^2]\) and \(t=\lfloor x^{401}\rfloor\), all groups \(\Sha_\ell(E_x,e)\) with \(1\le e<t\) are finite and \(V(E_x,t)\le\sqrt b\).
\end{theorem}

The auxiliary primes \(q\), \(d\) are fixed through the window while the field \(E_x\) changes with \(x\).  By Problem~\ref{mi-pr:split},
\[
 V(E_x,t)=B_2\log2+B_-+B_+,
\]
with
\[
 B_2=\sum_{1\le e<t}h_{2,e}(E_x),\qquad B_\pm=\sum_{\ell\ \mathrm{odd}}\log\ell\sum_{1\le e<t}h^\pm_{\ell,e}(E_x).
\]
Theorem~\ref{mi-thm:minus} bounds \(B_-\) at every scale.  Theorem~\ref{pl-thm:plus} bounds \(B_+\), given a bound on the layers at the bad primes.  The prime \(2\) and those layers are where the window is needed.

\begin{problem}\label{pr:page}
For \(x\ge1\), \(a(x)+2\le4(1+\log x)\).
\end{problem}

\begin{problem}\label{pr:log-qd}
If \((q,d)\) is a window pair at \(Y_0\), then \(\log(qd)\le Y_0^{\varepsilon/4}\) and, for each fixed \(J\ge1\), \(q,d>2Y_0^{401J}+1\) for all large \(Y_0\), uniformly in the pair.
\end{problem}

\section{The prime two and the bad primes}\label{sec:two}

At \(\ell=2\) the field \(E_x\) is compared with a layer over a base field that does not depend on \(x\).  Let \(L_0=A_\psi(i,\sqrt{-q},\sqrt{-d})\).

\begin{problem}\label{pr:L0}
\(i\in L_0\), \([L_0:\Q]\le8o\), \(L_0\subseteq\Q(\zeta_{4N_\psi qd})\), and \(E_x\subseteq L_0(\zeta_{2^{a(x)+2}})\).
\end{problem}

\begin{lemma}[descent at two]\label{lem:two}
For every \(a\ge2\) there is a layer \(L\) of degree \(2^m\) over \(L_0\), with \(m\le a+2\), such that \(h_{2,e}(E(\psi,a,q,d))\le1+h_{2,e}(L)\) for every \(e\ge1\).
\end{lemma}

\begin{lemma}[uniform towers]\label{lem:uniform}
There is \(R_2\), depending only on \(\psi\), such that for every pair of primes \(q\equiv3\pmod8\), \(d\equiv7\pmod{16}\) prime to \(N_\psi\) the hypothesis of Lemma~\ref{tw-lem:tower} holds over \(L_0\) with \(\ell=2\) and \(R=R_2\).  For each odd \(\ell\in\Bc_\psi\) there is \(R_\ell\) such that it holds over \(\Q(\zeta_{\ell N_\psi})\) with \(R=R_\ell\).
\end{lemma}

\noindent\emph{Status}: both lemmas are proved in the Lean formalization, the second from Theorem~\ref{tw-thm:fw}; we have read their statements only.  We have not followed how the formalization makes \(R_2\) independent of \(q\) and \(d\); the congruences on \(q\) and \(d\) are among its hypotheses.

\begin{problem}\label{pr:L0-disc}
Let \(j_0=\lfloor\log_2R_2\rfloor+3\) and \(L'\) the layer of degree \(2^{j_0}\) over \(L_0\).  Then \([L':\Q]\le8o\,2^{j_0}\), and \(\log(2\lvert\Delta_{L'}\rvert)\le c_1(1+\log(qd))\) with \(c_1\) depending only on \(\psi\).
\end{problem}

\section{The pigeonhole on scales}\label{sec:pigeonhole}

By Theorem~\ref{tw-thm:point}, for \(\ell=2\) over \(L_0\), and for each \(\ell\in\Bc_\psi\) over \(\Q(\zeta_{\ell N_\psi})\), the twists \(e\le Z\) split into at most \(R_\ell\) exceptional ones and the rest.  At the exceptional twists only Theorem~\ref{tw-thm:H2} is available, and by Problem~\ref{tw-pr:H2-small} it is affordable only for \(e\le x^\eta\), with \(\eta=\varepsilon/(2C_H)\).  The twists that matter at scale \(x\) are those with \(e\le x^{C_0}\).  So one needs scales at which no exceptional twist lies in \((x^\eta,x^{C_0}]\).

\begin{problem}\label{pr:window1}
Let \(R\ge0\), \(B\ge1\), \(U>0\), and \(T\) a set of at most \(R\) real numbers.  There is \(y\) with \(U/(2^{R+2}B^R)\le y\le U/2\) such that for every \(\alpha\in T\) and \(u\in[y,2y]\), \(\alpha\le u\) implies \(B\alpha\le u\).
\end{problem}

\begin{problem}\label{pr:window}
Let \(0<\eta\le C_0\), \(R\ge0\), \(X>1\), \(\Ec\) a set of at most \(R\) positive integers, and \(J=2^{R+2}(C_0/\eta)^R\).  There is \(Y\) with \(X^{1/J}\le Y\) and \(Y^2\le X\) such that for every \(x\in[Y,Y^2]\) and \(e\in\Ec\), \(e\le x^{C_0}\) implies \(e\le x^\eta\).
\end{problem}

\begin{problem}\label{pr:zhang}
In Problem~\ref{pr:window} let \(L=\log X\) and \(u_i=(L/2)(\eta/(2C_0))^i\) for \(0\le i\le R\).  The intervals \((\eta u_i,2C_0u_i]\) are disjoint, one of them contains no \(\log e\) with \(e\in\Ec\), and \(Y=e^{u_i}\) for that \(i\) satisfies the conclusion with \(J\) replaced by \(J/2\).
\end{problem}

\begin{problem}\label{pr:window-example}
For \(R=1\), \(\Ec=\{1000\}\), \(C_0=2\), \(\eta=1/2\) and \(X=e^{20}\), the proof of Problem~\ref{pr:window1} by induction on \(R\) gives \(Y=e^{5/4}\); here \(J=32\), and \(x^2\le e^5<1000\) for every \(x\in[Y,Y^2]\).
\end{problem}

\begin{problem}\label{pr:union}
One window serves all the primes at once: if \(\Ec_2\) and \(\Ec_\ell\), \(\ell\in\Bc_\psi\), have at most \(R_2\) and \(R_\ell\) elements, Problem~\ref{pr:window} applies to their union with \(R=R_2+\sum_{\ell\in\Bc_\psi}R_\ell\).  With \(\eta=\varepsilon/(2C_H)\) the exponent is \(J=2^{R+2}(2C_HC_0/\varepsilon)^R\).
\end{problem}

\section{The bounds on a window}\label{sec:bounds}

Fix \(C_0\ge1\) and a function \(t\) with \(2\le t(x)\le x^{C_0}\) for \(x\ge2\).  Let \(X\) be large, \(q\equiv3\pmod8\) and \(d\equiv7\pmod{16}\) primes prime to \(N_\psi\) with \(\log(qd)\le X^{\varepsilon/(4J)}\), \(Z=\lfloor X^{C_0}\rfloor\), and \(Y\) as in Problem~\ref{pr:union} for the exceptional sets of Theorem~\ref{tw-thm:point} with this \(Z\).  Let \(x\in[Y,Y^2]\).

\begin{problem}\label{pr:nonexc}
Let \(e\le t(x)\) be non-exceptional for \(\ell=2\), and \(L\) a layer of degree \(2^m\) over \(L_0\) with \(m\le a(x)+2\).  Then \(h_{2,e}(L)\le C(1+\log x)\), with \(C\) independent of \(x\), \(X\), \(q\) and \(d\).
\end{problem}

\begin{problem}\label{pr:exc}
Let \(e\le t(x)\) be exceptional for \(\ell=2\), and \(L\) as in Problem~\ref{pr:nonexc}.  Then \(h_{2,e}(L)\le C(1+\log x)+2\exp(Cx^{3\varepsilon/4})\).
\end{problem}

\begin{problem}\label{pr:final}
For every \(C\) and all large \(x\), \(x^{C_0}\bigl(1+C(1+\log x)+2\exp(Cx^{3\varepsilon/4})\bigr)\le\exp(x^{4\varepsilon/5})\).
\end{problem}

The same bounds hold, more easily, at \(\ell\in\Bc_\psi\) over the fixed layer of degree \(\ell^{j_\ell}\), \(j_\ell=\lfloor\log_\ell R_\ell\rfloor+3\), over \(\Q(\zeta_{\ell N_\psi})\), whose discriminant does not depend on \(q\) and \(d\).  Problems~\ref{pr:page}, \ref{pr:L0-disc} and \ref{pr:window1}--\ref{pr:final}, Lemmas~\ref{lem:two} and~\ref{lem:uniform}, and Theorem~\ref{tw-thm:point} prove the following.

\begin{theorem}[windows at two and at the bad primes]\label{thm:bab}
Let \(C_0\ge1\), \(A_1>0\), \(R=R_2+\sum_{\ell\in\Bc_\psi}R_\ell\) and \(J=2^{R+2}(2C_HC_0/\varepsilon)^R\).  There is \(X_0\) such that for all functions \(a\), \(t\) with \(2\le a(x)\), \(2^{a(x)}\le x^{A_1}\) and \(2\le t(x)\le x^{C_0}\) for \(x\ge2\), every \(X\ge X_0\) and all primes \(q\equiv3\pmod8\), \(d\equiv7\pmod{16}\) prime to \(N_\psi\) with \(\log(qd)\le X^{\varepsilon/(4J)}\), there is \(Y\) with \(X^{1/J}\le Y\) and \(Y^2\le X\) such that for every \(x\in[Y,Y^2]\)
\[
 \sum_{1\le e<t(x)}h_{2,e}(E_x)\le\exp(x^{4\varepsilon/5}),
\]
and the same bound holds for \(\sum_{1\le e<t(x)}h_{\ell,e}\) of the layer of degree \(\ell^{j_\ell}\) over \(\Q(\zeta_{\ell N_\psi})\), for every \(\ell\in\Bc_\psi\).
\end{theorem}

\noindent\emph{Status}: proved in the Lean formalization, as just described.  Problem~\ref{pr:page} uses the page exponent; for general \(a\) the bound \(m\le a+2\le C(1+\log x)\) follows from \(2^{a(x)}\le x^{A_1}\) in the same way.  Zhang states the result as \cite[Prop.~5.3]{Zh}, with the window lemma \cite[Lemma~5.5]{Zh}.

\section{The assembly}\label{sec:assembly}

Let \(t=\lfloor x^{401}\rfloor\), and let \(x\) lie in the window of Theorem~\ref{thm:bab} with \(C_0=401\) and \(A_1=4\); these are admissible since \(2x^{2+8\varepsilon}\le x^4\) for \(x\ge2\).  By Problem~\ref{lt-pr:conductor}, \(\log(N_xt)\le2x^{\varepsilon/4}\) for large \(x\).

\begin{problem}\label{pr:minus-bound}
The right side of Theorem~\ref{mi-thm:minus} for \(E_x\) and \(t\) is at most \(\exp(Cx^{\varepsilon/4})\), at every large scale \(x\).
\end{problem}

\begin{problem}\label{pr:plus-bound}
Theorem~\ref{pl-thm:plus} applies on the window, with \(B=\exp(x^{4\varepsilon/5})\), and gives \(B_+\le\exp(2Cx^{\varepsilon/4})+C\exp(x^{4\varepsilon/5})\).
\end{problem}

\begin{problem}\label{pr:total}
On the window, all \(\Sha_\ell(E_x,e)\), \(1\le e<t\), are finite and \(V(E_x,t)\le\sqrt b\) for all large \(x\).
\end{problem}

Problems~\ref{pr:log-qd}, \ref{pr:minus-bound}--\ref{pr:total} and Theorem~\ref{thm:bab} prove Theorem~\ref{thm:vwin}.  \emph{Status}: this is how the Lean formalization proves Theorem~\ref{thm:vwin}; the inputs it uses are those listed in Chapters~\ref{tw-chap}, \ref{mi-chap} and~\ref{pl-chap}, with their statuses there.  In \cite{Adr} the statement is assumed as a hypothesis, and Chapter~\ref{lt-chap} records it as such.

\paragraph{Why windows.}
The minus part and the plus part away from \(\Bc_\psi\) are bounded at every scale.  At \(\ell=2\) and at the finitely many bad primes, the vanishing of \(\mu\) makes the lengths along a cyclotomic tower linear in the layer, except at boundedly many twists \(\ell\)-adically close to fixed centres (Theorem~\ref{tw-thm:point}).  At those twists only the doubly exponential bound for \(H^2\) is known, which is affordable only for \(e\le x^{\varepsilon/(2C_H)}\) (Problems~\ref{tw-pr:H2-small} and~\ref{tw-pr:H2-large}).  The exceptional twists do not move with \(x\), so the scale is moved instead: a pigeonhole on logarithmic scales finds \([Y,Y^2]\) free of them, at the price of the exponent \(J\), which grows like \((C/\varepsilon)^R\).  Chapter~\ref{lt-chap} then passes from \([Y,Y^2]\) to ranges \([Y,Y^{3/2}]\) (Theorem~\ref{lt-thm:W3}); in the formalization a similar passage, to families of scales with parameter \(3/2\), is made by a step we have not followed.

\begin{leanremark}
\raggedright In \cite{Lean} Theorem~\ref{thm:vwin} is \tl{VWin} (\tl{VWinAt}~\(\psi\)~\(401\) for every \(\psi\)), with \tl{VAt}, \tl{WinPair}, \tl{PageN} and \tl{IsChiField}; it is proved by \tl{UpperTheory.\allowbreak{}One.\allowbreak{}vwin\_\allowbreak{}of\_\allowbreak{}row} through \tl{HV.\allowbreak{}Core.\allowbreak{}of\_\allowbreak{}four\_\allowbreak{}alt}, \tl{vwin\_\allowbreak{}of\_\allowbreak{}counts} and \tl{vwinAt\_\allowbreak{}of\_\allowbreak{}joint}, whose last steps are \tl{vvle\_\allowbreak{}of\_\allowbreak{}parts}, \tl{admissible\_\allowbreak{}of\_\allowbreak{}pageN}, \tl{log\_\allowbreak{}mul\_\allowbreak{}le\_\allowbreak{}rpow\_\allowbreak{}of\_\allowbreak{}winPair}, \tl{eventually\_\allowbreak{}winPair\_\allowbreak{}gt\_\allowbreak{}power} and \tl{error\_\allowbreak{}terms\_\allowbreak{}le\_\allowbreak{}sqrt\_\allowbreak{}exp\_\allowbreak{}rpow}.  Theorem~\ref{thm:bab} is \tl{ThmBab}, proved by \tl{Win.\allowbreak{}thmBab\_\allowbreak{}of} from \tl{partA} (\(\ell=2\)) and \tl{partB} (\(\ell\in\Bc_\psi\)); Problems~\ref{pr:window1} and~\ref{pr:window} are \tl{Win.\allowbreak{}window1} and \tl{Win.\allowbreak{}window}, the conclusion being \tl{Good}; Problem~\ref{pr:final} is \tl{Win.\allowbreak{}final}.  Lemma~\ref{lem:two} is \tl{TwoBase}, and Lemma~\ref{lem:uniform} is \tl{TwoIn} and \tl{FixedIn}, from \tl{towerTwo} and \tl{towerFixed}.  The bound is passed on by \tl{supplyChi}, as \tl{HostExW}~\(3/2\).
\end{leanremark}

\section*{Hints}

\subsection*{\S\ref{sec:statement}}
\begin{description}
\item[\ref{pr:page}] \(\varepsilon<1/52\) gives \((2+8\varepsilon)/\log2<3.2\).
\item[\ref{pr:log-qd}] The size conditions in the definition of a window pair; \(Y_0^{\varepsilon/20}\) eventually exceeds \(2(1+\log^2Y_0)\).
\end{description}

\subsection*{\S\ref{sec:two}}
\begin{description}
\item[\ref{pr:L0}] \(q\equiv d\equiv3\pmod4\), so \(\sqrt{-q}\in\Q(\zeta_q)\) and \(\sqrt{-d}\in\Q(\zeta_d)\); compare the four factors of \(E_x\).
\item[\ref{pr:L0-disc}] Problem~\ref{tw-pr:layer-disc} with \(N_0=4N_\psi\) and \(P=qd\).
\end{description}

\subsection*{\S\ref{sec:pigeonhole}}
\begin{description}
\item[\ref{pr:window1}] Induction on \(R\).
\item[\ref{pr:window}] Problem~\ref{pr:window1}, with \(u=\log x\).
\item[\ref{pr:zhang}] \(2C_0u_{i+1}=\eta u_i\).
\item[\ref{pr:union}] The cardinality of a union.
\end{description}

\subsection*{\S\ref{sec:bounds}}
\begin{description}
\item[\ref{pr:nonexc}] Theorem~\ref{tw-thm:point} with \(\ell=2\) and Problem~\ref{pr:page}; \(\log Z\le C_0J\log x\), and \(2^{c_0}\le16o\) by Problem~\ref{tw-pr:root-degree}.
\item[\ref{pr:exc}] Theorem~\ref{tw-thm:point}, Problems~\ref{pr:L0-disc} and~\ref{tw-pr:H2-small}; the window gives \(e\le x^{\varepsilon/(2C_H)}\).
\item[\ref{pr:final}] \(4/5>3/4\).
\end{description}

\subsection*{\S\ref{sec:assembly}}
\begin{description}
\item[\ref{pr:minus-bound}] \([E_x:\Q]\le N_x\).
\item[\ref{pr:plus-bound}] Problem~\ref{pr:log-qd} for the size conditions on \(q\) and \(d\); Theorem~\ref{thm:bab} for the layers.
\item[\ref{pr:total}] Theorem~\ref{thm:bab} with Lemma~\ref{lem:two}, and Problems~\ref{pr:minus-bound}, \ref{pr:plus-bound}; all exponents are below \(\varepsilon\).
\end{description}

\begin{thebibliography}{Lean}
\bibitem[Adr]{Adr} J.~Adriazola, \emph{The Alp\"oge proof of the Riemann hypothesis}, manuscript, revised 24 November 2026, \url{https://github.com/jimjamzola/RH-Alpoge-Reconstruction}.
\bibitem[Lean]{Lean} The Lean community, \emph{riemann-hypothesis}, a Lean formalization of Alp\"oge's argument, \url{https://github.com/leanprover-community/riemann-hypothesis}.
\bibitem[Zh]{Zh} R.~Zhang, \emph{A preliminary note on Alp\"oge's proposed proof of the generalized Riemann hypothesis}, 22 November 2026, \url{https://hexagonmath.org/pdf/2611.00809v1}.
\end{thebibliography}

\end{document}
